\documentclass[10pt,a4paper]{amsart}
\usepackage[margin=19mm]{geometry}
\usepackage{amsmath,amssymb,mathtools}
\usepackage[scr=rsfs,cal=euler]{mathalfa}
\usepackage{xfrac}
\usepackage[unicode,hidelinks,bookmarks=false]{hyperref}
\newcommand{\ie}{i.e.\ }
\newcommand{\cf}{cf.\ }
\newcommand{\resp}{respectively\ }
\newcommand{\Subsection}{\subsection}
\providecommand{\nobreakdash}{\nobreak\hbox{-}}
\setlength{\emergencystretch}{2em}
\multlinegap0pt
\usepackage{bm}
\usepackage{bbm}
\usepackage{dsfont}
\usepackage{xspace}
\usepackage{cancel}
\usepackage{mathtools}
\usepackage{amsthm}
\makeatletter
\@ifundefined{thm}{\newtheorem{thm}{Thm}}{}
\@ifundefined{lemma}{\newtheorem{lemma}[thm]{Lemma}}{}
\makeatother
\makeatletter
\expandafter\ifx\csname algo\endcsname\relax
\newenvironment{algo}{\begin{algorithm}\rm}{\end{algorithm}}
\fi
\expandafter\ifx\csname assum\endcsname\relax
\newenvironment{assum}{\begin{assumption}\rm}{\end{assumption}}
\fi
\expandafter\ifx\csname clai\endcsname\relax
\newenvironment{clai}{\begin{claim}\rm}{\end{claim}}
\fi
\expandafter\ifx\csname de\endcsname\relax
\newenvironment{de}{\begin{definition}\rm}{\end{definition}}
\fi
\expandafter\ifx\csname ex\endcsname\relax
\newenvironment{ex}{\begin{example}\rm}{\end{example}}
\fi
\expandafter\ifx\csname exmp\endcsname\relax
\newenvironment{exmp}{\begin{example}\rm}{\end{example}}
\fi
\expandafter\ifx\csname fa\endcsname\relax
\newenvironment{fa}{\begin{fact}\rm}{\end{fact}}
\fi
\expandafter\ifx\csname ob\endcsname\relax
\newenvironment{ob}{\begin{observation}\rm}{\end{observation}}
\fi
\expandafter\ifx\csname prob\endcsname\relax
\newenvironment{prob}{\begin{problem}\rm}{\end{problem}}
\fi
\expandafter\ifx\csname qustion\endcsname\relax
\newenvironment{qustion}{\begin{question}\rm}{\end{question}}
\fi
\expandafter\ifx\csname re\endcsname\relax
\newenvironment{re}{\begin{remark}\rm}{\end{remark}}
\fi
\expandafter\ifx\csname rem\endcsname\relax
\newenvironment{rem}{\begin{remark}\rm}{\end{remark}}
\fi
\makeatother
\title[On the future cover of a sofic shift]{On the future cover of a sofic shift}
\author{Klaus Thomsen}
\date{Source: arXiv:2512.01368v3 (CC BY 4.0)}
\begin{document}
\maketitle

\noindent\emph{Source and licence.} On the future cover of a sofic shift. Version arXiv:2512.01368v3 (CC BY 4.0), distributed under the Creative Commons Attribution 4.0 International licence, as recorded in the arXiv licence metadata for this exact version and shown on the abstract page https://arxiv.org/abs/2512.01368v3. Full rights notice: licenses/arxiv-2512.01368-CC-BY-4.0.txt.\par\smallskip

\noindent\textit{Editorial scope.} Counted unit: the Lemma whose source label is \texttt{02-09-25dx} in the article (source0.tex lines 875--880, source label \texttt{02-09-25dx}), delivered together with the article's own complete original proof of it (source0.tex lines 881--922, one \texttt{proof} environment of 42 lines). No mathematics is written, repaired or re-derived: statement and proof are the article's own text line for line. Required context retained uncounted: 02-09-25ax (lines 612--615); 11-09-25g (lines 808--810). Every definition, display and label of those lines is retained unchanged. Omitted from this package and not redistributed: 1--611, 616--807, 811--874, 923--3937. Nothing inside a retained excerpt is omitted or altered: every retained body is delivered exactly as the article's lines are written, so no omission receipt of any kind accompanies this package. Citations: the retained excerpts carry no citation command at all, so no bibliography entry of the article is transcribed and no reference is redistributed. Reference adaptation: none was needed and none was made; every reference inside the delivered unit resolves inside it, so no unresolved reference or citation is produced. Printed slips kept literal: no printed word, symbol, hypothesis or number of the counted statement or of its proof was altered. Layout and class: the article's own class is \texttt{amsart} with packages including tcolorbox, graphicx, caption, amsmath, amssymb, hyperref, enumerate, color, xcolor, mathtools, float, tikz, bm, appendix; the wrapper is the lane's pinned \texttt{amsart} editorial wrapper at 10pt on A4, which restates the theorem-like environments the retained text needs and adds the article's own macro definitions that the retained text needs (none), with extra wrapper packages (bm, bbm, dsfont, xspace, cancel, mathtools). The delivered numbering is the unit's own. Assets: no retained span contains a figure environment, a graphics-inclusion command, a PDF-page inclusion, a table or a TikZ picture; no image file is copied into this package, the package declares no asset dependency and no image file is redistributed with it. Licence: version arXiv:2512.01368v3 of this article is distributed under the Creative Commons Attribution 4.0 International licence, as recorded in the arXiv licence metadata for this exact version and shown on the abstract page https://arxiv.org/abs/2512.01368v3. A full rights notice is delivered at licenses/arxiv-2512.01368-CC-BY-4.0.txt. Characters: every retained excerpt is pure ASCII, so the delivered engine, XeLaTeX with its default Latin Modern text font, reports no missing character.
\begin{equation}\label{02-09-25ax}
\left\{ L_H(\mu) : \ \mu \in X_H[0,\infty), \ s_H(\mu)  = 
t_H(x_{(-\infty, k]})\right\} = \left\{ y \in Y[0,\infty): \ L_H(x_{(-\infty,k]})y \in Y \right\}, 
\end{equation}

\begin{lemma}\label{11-09-25g} $\alpha_Y(y)$ is regular in $X_{\mathbb K(Y)}$ for all $y \in Y$, i.e.
$\alpha_Y(Y) \subseteq \mathcal R(L_{\mathbb K(Y)})$.
\end{lemma}

\begin{lemma}\label{02-09-25dx} Assume that $(G,L_G)$ is a right-resolving and regular presentation of $Y$. It follows that there is a labeled-graph homomorphism $\theta : (G,L_G) \to (\mathbb K(Y), L_{\mathbb K(Y)})$ such that
$$
\theta\left(\mathcal R(L_G)\right) \subseteq \mathcal R(L_{\mathbb K(Y)}).
$$
If $(G,L_G)$ also is follower-separated, $\theta : G \to \mathbb K(Y)$ is injective.
\end{lemma}

\begin{proof} Let $v \in V_G$. Since $v$ is regular there is an $x \in X_G$ such that $t_G(x_{(-\infty,-1]}) = v$ and 
$f_G(v)  = F(L_G(x))$.
Define $\theta(v) \in V_{\mathbb K(Y)}$ such that 
$$
\theta(v) := F(L_G(x)) .
$$
This is well-defined: If $x'$ has the same properties as $x$, it follows that
$F(L_G(x')) = f_G(v) = F(L_G(x))$.


Let $e \in E_G$. Choose $x \in X_G$ such that $t_G(x_{(-\infty,-1]}) = s_G(e)$ and 
$f_G(s_G(e)) = F(L_G(x))$.
Set $x'_i = x_{i+1}, \ i \leq -2$, $x'_{-1} =e$ and let $x'_i, i \geq 0$, be arbitrary subject to the condition that $x' \in X_G$. Then $t_G(x'_{(-\infty, -1]}) = t_G(e)$ and by using that $(G,L_G)$ is right-resolving it follows that
\begin{align*}
& f_G(t_G(e)) = \frac{f_G(s_G(e))}{L_G(e)}
 = \frac{\left\{y \in Y[0,\infty): \ L_G(x_{(-\infty,-1]})y \in Y \right\}}{L_G(e)} \\
 &= \left\{y \in Y[0,\infty): \ L_G(x'_{(-\infty,-1]})y \in Y \right\}.
\end{align*}
Hence $\theta(t_G(e)) = F(L_G(x'))$ and  $F(L_G(x')) = \frac{F(L_G(x))}{L_G(e)}$. There is therefore an edge $e'$ in $\mathbb K(Y)$ from $\theta(s_G(e)) =F(L_G(x)) $ to $\theta(t_G(e)) = F(L_G(x'))$ with label $L_{\mathbb K(Y)}(e') =L_G(e)$. Since $(\mathbb K(Y), L_{\mathbb K(Y)})$ is right-resolving this arrow is unique and we set $\theta(e):= e'$. Then $\theta : G \to \mathbb K(Y)$ is a labeled-graph morphism.

When $x \in \mathcal R(L_G)$, it follows from the definition of $\theta$ and the defining relation \eqref{02-09-25ax} for regular rays that $\theta(x) \in X_{\mathbb K(Y)}$ is the path
$$
\cdots  \overset{L_G(x_{k-1})}{\to}F(L_G(\sigma^k(x))) \overset{L_G(x_k)}{\to} F(L_G(\sigma^{k+1}(x))) \cdots .
$$ 
Thus $\theta(x) = \alpha_Y(L_G(x))$ and hence $\theta(x) \in \mathcal R(L_{\mathbb K(Y)})$ by Lemma \ref{11-09-25g}. 






%Since, by Lemma \ref{02-10-24fx},
%\begin{align*}
%&f_{\mathbb K(Y)}(t_{\mathbb K(Y)}(\theta(x)_{(-\infty, k]})) = f_{\mathbb K(Y)}(F(L_G(\sigma^{k+1}(x)))) = F(L_G(\sigma^{k+1}(x)) \\
%&= \left\{ y \in Y[0,\infty): \ L_G(x_{(-\infty, k]})y \in Y \right\} \\
%& = \left\{ y \in Y[0,\infty): \ L_{\mathbb K(Y)}(\theta(x)_{(-\infty, k]})y \in Y \right\}  \\
%& = F(\sigma^{k+1}(L_{\mathbb K(Y)}(\theta(x))))
%\end{align*}
%for all $k \in \mathbb Z$, we see that $\theta(x) \in \mathcal R(L_{\mathbb K(Y)})$, cf. \eqref{23-09-25}.

Finally, assume that $v,w \in V_G$ are such that $\theta(v) =\theta(w)$. There are elements $x,x' \in X_G$ such that $t_G(x_{(-\infty,-1]}) = v$ and $t_G(x'_{(-\infty,-1]}) = w$, $f_{G}(v) = F(L_G(x))$ and $f_G(w) = F(L_G(x'))$. By definition of $\theta$ this implies that $f_G(v) = \theta(v) = \theta(w) = f_G(w)$ and hence that $v=w$ when $(G,L_G)$ is follower-separated.
\end{proof}

\end{document}
